 \section{Introduction}
 
 Wave equations with memory effects have been a central topic in the analysis of partial
 differential equations, owing to their applications in physics, engineering, and material
 science. The focus has been on the existence, stability, and energy decay of solutions, with significant progress made for both single and coupled wave equations.
 To motivate our work, data collected and results are reviewed to provide a thorough study of
 both the single and coupled viscoelastic wave equation. We start with the pioneer works of Dafermos \cite{Dafermos1, Dafermos2}, where the author showed that solutions decay to zero
 for smooth, monotone relaxation functions $g$, but the decay rate remained unspecified. Cavalcanti et al.\ \cite{CavalcantiDS} studied
\[
|u_t|^\rho u_{tt} - \Delta u - \Delta u_{tt} + \int_0^t g(t - s) \Delta u(s) \, ds - \gamma \Delta u_t = 0,\quad \text{in } \Omega \times \mathbb{R}^{*}_{+}
\]
proving global existence for $\gamma \geq 0$ and exponential energy decay for $\gamma > 0$. Messaoudi and Tatar \cite{MTatar1, MTatar2} extended these results to include source terms, showing exponential and polynomial decay depending on the rate of decay of g(t).
Cavalcanti et al.\ \cite{CavalcantiFerreira} considered a viscoelastic wave equation with a
local frictional damping, where the kernel $g$ satisfies, for two positive constants $\xi_1$
and $\xi_2$,
\begin{align*}
-\xi_1 g(t) \leq g'(t) \leq -\xi_2 g(t), \quad \forall t \in \mathbb{R}_{+}.
\end{align*}
Under the above assumptions and subject to certain restrictions on the control zone, they established an exponential decay result. Berrimi and Messaoudi \cite{Berrimi2004} later
improved this result by proving that the viscoelastic dissipation alone is sufficient to stabilize the system, and subsequently extended it to the case where a source term competes with the dissipation \cite{Berrimi2006}.
Further contributions including the analysis of nonlinear damping with polynomially decaying kernels \cite{Cavalcanti2003} and global existence with uniform decay for source-type problems \cite{Li2011}  appeared later.
 Without restricting $g(t)$ to exponential or polynomial decay, Messaoudi \cite{SMessaoudi1, SMessaoudi2} generalized all the existing decay results by considering the following equation
\[
u_{tt} - \Delta u + \int_0^t g(t - \tau) \Delta u(\tau) \, d\tau = b |u|^{\gamma} u,\quad b \geq 0,
\]
where $g'(t) \leq -\xi(t) g(t)$, for a non-increasing differentiable $\xi(t)$.
These latter results inspired other authors and more results have been established, see \cite{SaidH2011, Hao, PJS, WSH, MessaoudiMustafa}.

Further improvement was establish by Messaoudi and Al-Khulaifi \cite{MessaoudiWaled}, who addressed the problem in \cite{CavalcantiDS}, for $\gamma =0$ and for a kernel satisfying, for $1\leq p\leq \frac{3}{2}$, the condition
\[
g'(t) \leq -\xi(t) g^p(t), \quad \forall t \in \mathbb{R}_{+}.
\]
They proved that the energy of solutions converges to zero at the same rate as $g(t)$, provided $g(t)$ decays faster than $s^{-2}$ at infinity. Later, Mustafa \cite{MM181} extended his earlier joint work with Messaoudi \cite{MM}, which dealt with kernels satisfying $g'(t) \leq -G(g(t))$, to the more general condition $$g'(t) \leq -\xi(t) G(g(t)), \quad \forall t \in \mathbb{R}_{+},$$ where $G$ is strictly increasing and convex function. This extension provided explicit energy decay estimates and encompassed a wider range of decay behaviors.

Coupled wave systems also have consistently drawn attention due to their fundamental role in understanding stability and long-term dynamics. In a recent contribution, Guesmia, Messaoudi, and Zahri \cite{GASM} analyzed an abstract coupled system involving a viscoelastic component interacting with a purely elastic one. Precisely, they looked into the problem
\begin{gather*}
u_{tt}(t) + Au(t) - \int_0^{t} g(s) A^{\theta} u(t - s) \, ds + \alpha u + \beta B v(t) = f(u), \quad t > 0, \\
v_{tt}(t) + A v(t) + \beta B u(t) = 0, \quad t > 0,\\
u(0) = u_0, \quad u_t(0) = u_1, \quad v(0) = v_0, \quad v_t(0) = v_1,
\end{gather*}
 where $A$ is a self-adjoint positive operator on a Hilbert space $H$, $\theta \in [0,1]$,
 $\alpha > 0$, and $g$ is a positive nonincreasing relaxation function.
 Their analysis established general decay rates for the system energy, showing that the
 asymptotic behavior depends strongly on the relaxation function and the coupling parameter.
 Numerical experiments were also provided, highlighting the impact of memory effects on
 the qualitative behavior of solutions. Earlier, Feng and Li \cite{FengBaoweiHaiyan} considered
 a viscoelastic Lam\'e system with strong damping terms. For two distinct relaxation kernels
 $g_1$ and $g_2$, they assumed conditions of the form $g_i'(t) \leq -\xi_i(t) H_i(g_i(t))$
 ($i=1,2$), which led to explicit and general decay results. Their work improved earlier results
 of Beniani, Taouaf, and Benaissa \cite{Beniani}, who had previously established a general
 decay property for similar coupled viscoelastic structures.
Further investigations were carried out by Jin et al. \cite{JKJT}, who studied a coupled system of two abstract evolution equations with finite memory, given by
\begin{gather*}
 u_{tt}(t) + Au(t) - \int_0^{t} g(s) A u(t - s) \, ds + \alpha u + \beta B v(t)
 = f(u), \quad t > 0, \\
v_{tt}(t) + A v(t) + \beta B u(t) = 0, \quad t > 0,
 \end{gather*}
 where $A$ and $B$ are operators, $\alpha \geq 0$, $\beta > 0$, and $f$ is a nonlinear source.
 Under appropriate assumptions on $g$ and structural conditions on $A$, $B$, and $f$,
 they proved polynomial stability with decay rates of order $t^{-1}$. Several extensions and
related contributions in this direction may be found in \cite{ ORHC, AkilM}.
In addition to viscoelastic systems, purely elastic coupled wave models have also been studied.
For instance, a weakly coupled system
\begin{gather*}
u_{tt} - \Delta u + u_t + \alpha v = 0, \quad \text{in } \Omega \times (0, \infty), \\
v_{tt} - \Delta v + \alpha u = 0, \quad \text{in } \Omega \times (0, \infty),
\end{gather*}
was analyzed in \cite{LRF}, where $\Omega \subset \mathbb{R}^n$ is a bounded domain and
$\alpha$ a positive constant. The authors showed that the energy decays polynomially,
with the rate being optimal, and provided numerical validation in the one-dimensional case.
Similar results on stability and decay rates had earlier been obtained for systems with
boundary damping, notably by Komornik et al. \cite{KVBR} and Aassila \cite{Assila1, Assila2},
where exponential stability was achieved in the case of linear damping and polynomial stability
in the case of nonlinear damping mechanisms.

Motivated by the findings in \cite{GASM}, the present work investigates the coupled viscoelastic wave system
\begin{equation} \label{e1.1}
\begin{gathered}
u_{tt} - \Delta u + \int_0^t g(t - s) \Delta u \, ds + \alpha v = 0, \quad
 \text{in } \Omega \times (0, T), \\
v_{tt} - \Delta v + \alpha u = 0, \quad \text{in } \Omega \times (0, T), \\
u = v = 0, \quad \text{on } \partial \Omega \times (0, T), \\
u(0) = u_0, \, u_t(0) = u_1, \; v(0) = v_0, \, v_t(0) = v_1, \quad \text{in } \Omega,
\end{gathered}
\end{equation}
where \( \Omega \) is a bounded domain in \( \mathbb{R}^N, \) \( \alpha > 0 \), and the initial
data belong to suitable spaces. The relaxation function \( g \) satisfies
\( g'(t) \leq -\xi(t) G(g(t)) \), where $G$ is an increasing and convex function near the origin and
$\xi$ is a nonincreasing function. We establish, in details, an existence and uniqueness result,
then we prove a general and explicit result for energy decay and provide some numerical
illustration to validate our theoretical findings.

The structure of this paper is as follows: the next section presents some assumptions,
remarks and functionals needed in our results. In Section 3, we establish the existence of
solutions by combining energy estimates with the Faedo-Galerkin method.
Section 4 presents several technical lemmas and contains the statement and proof of the
decay results.
In section 5, two numerical tests are conducted aiming to verify our decay result.

 \section{Preliminaries}

 This section presents the foundation for our main result. We set
 \begin{gather*}
 (h \star \varphi)(t) = \int_0^t h(t-s) \varphi(s)ds,\
 (h \diamond \varphi)(t) = \int_0^t h(t-s)\| \varphi(t)- \varphi(s)\| ds, \\
 (h \circ \varphi)(t) = \int_0^t h(t-s)\int_\Omega\| \varphi(t)- \varphi(s)\|^2 dx \, ds.
 \end{gather*}
 For the relaxation function $g$, we assume:
 \begin{itemize}
\item[(A1)] $g: \mathbb{R}^+ \to \mathbb{R}^+ $ is a $ C^1$ function satisfying\\
\begin{equation*}
 g(0) > 0, \quad 1-\int_0^{\infty} g(s) \, ds = l > 0. \label{e2.1}
\end{equation*}

\item[(A2)] There exists a $C^1$ function $G :\mathbb{R}^+ \to \mathbb{R}^+ $
which is linear or it is strictly increasing and strictly convex $C^2$ function on
$(0, r]$, $r \leq g(0)$, with $G(0) = G'(0) = 0$, such that
\begin{equation} \label{e2.2}
g'(t) \leq -\xi(t) G(g(t)),\quad \forall t \geq 0,
\end{equation}
where $\xi$ is a positive non-increasing differentiable function.

\item[(A3)] The constant $\alpha $ is such that
\begin{equation} \label{e2.3}
0 < \alpha < \frac{\sqrt{l} }{C_p},
\end{equation}
where $C_p$ is the Poincar\'e constant.
\end{itemize}

\begin{lemma}[\cite{Munoz}] \label{lem1}
For any function $h \in C^1 (\mathbb{R})$ and any $k \in H^1 (0,T)$, we have
$$
(h \star k)(t) k_t
 = -\frac{1}{2} h(t)  \| k\|^2 + \frac{1}{2} (h' \circ k)(t)
 -\frac{1}{2} \frac{d}{dt} \Big\{( h \circ k)(t)
 - \Big( \int_0^t h(\tau) d\tau \Big) \|k\|^2 \Big\}.
$$
\end{lemma}

\begin{lemma}[\cite{SMessaoudi1}] \label{lem2}
 For $u \in H^1_0(\Omega)$, we have
\begin{equation}
\int_\Omega \Big( \int_0^t g(t - s)(u(t) - u(s)) \, ds \Big)^2 dx
\leq (1 - l) C_p^2 (g \circ \nabla u)(t),\label{e2.4}
\end{equation}
where $l$ is given in {\rm (A1)}.
\end{lemma}

\section{Existence of weak solutions}
Now, we examine the well-posedness of problem \eqref{e1.1}. To establish this, we will combine some energy estimates and the Faedo-Galerkin method.

